% ==============================================================================
% Figura: Máquina de Estados Finitos (FSM) del Firmware furuta_control.c (1 kHz)
% Diseño por construcción (Layout by Construction): posicionamiento relativo,
% anclas exactas, contenedores elásticos (fit) y cero cruces de transiciones.
% Bibliotecas TikZ requeridas: arrows.meta, positioning, fit, calc, babel, bending.
% ==============================================================================
\begin{figure*}[!t]
\centering
% --- Paleta semántica suave (ámbito local al entorno figure*) ---
\definecolor{fsmIdleFill}{HTML}{F1F5F9}%
\definecolor{fsmIdleBdr}{HTML}{475569}%
\definecolor{fsmIdleHdr}{HTML}{334155}%
\definecolor{fsmServFill}{HTML}{EFF6FF}%
\definecolor{fsmServBdr}{HTML}{2563EB}%
\definecolor{fsmServHdr}{HTML}{1D4ED8}%
\definecolor{fsmServZoneBdr}{HTML}{60A5FA}%
\definecolor{fsmServZonePill}{HTML}{DBEAFE}%
\definecolor{fsmServZoneTxt}{HTML}{1E3A8A}%
\definecolor{fsmWarmFill}{HTML}{FFF7ED}%
\definecolor{fsmWarmBdr}{HTML}{EA580C}%
\definecolor{fsmWarmHdr}{HTML}{C2410C}%
\definecolor{fsmHeroFill}{HTML}{DCFCE7}%
\definecolor{fsmHeroBdr}{HTML}{059669}%
\definecolor{fsmHeroHdr}{HTML}{065F46}%
\definecolor{fsmAutoZoneBdr}{HTML}{0D9488}%
\definecolor{fsmAutoZonePill}{HTML}{CCFBF1}%
\definecolor{fsmAutoZoneTxt}{HTML}{115E59}%
\definecolor{fsmStopFill}{HTML}{FEF2F2}%
\definecolor{fsmStopBdr}{HTML}{DC2626}%
\definecolor{fsmStopHdr}{HTML}{B91C1C}%
\definecolor{fsmAlertFill}{HTML}{FFF1F2}%
\definecolor{fsmAlertBdr}{HTML}{E11D48}%
\definecolor{fsmAlertTxt}{HTML}{9F1239}%
\definecolor{fsmInk}{HTML}{1E293B}%
\definecolor{fsmMuted}{HTML}{475569}%
% --- Parámetros centralizados de geometría (Single Source of Truth) ---
\def\fsmNodePadX{7pt}%
\def\fsmColGapL{3.10cm}%
\def\fsmColGapR{3.35cm}%
\def\fsmIdleW{2.60cm}%
\def\fsmIdleH{2.20cm}%
\def\fsmServW{3.55cm}%
\def\fsmServH{1.00cm}%
\def\fsmServGapY{0.78cm}%
\def\fsmAutoW{2.85cm}%
\def\fsmAutoH{1.15cm}%
\def\fsmAutoGapX{2.05cm}%
\def\fsmAutoGapY{1.35cm}%
\def\fsmLaneStep{0.38cm}%
\resizebox{0.88\textwidth}{!}{%
\begin{tikzpicture}[
    font=\sffamily,
    % --- Estilos de nodos de estado (Layout by Construction) ---
    fsmNode/.style={
        rectangle,
        rounded corners=7pt,
        align=center,
        inner xsep=\fsmNodePadX,
        inner ysep=6pt,
        text=fsmInk,
        line width=1.05pt
    },
    stateIdle/.style={
        fsmNode,
        fill=fsmIdleFill,
        draw=fsmIdleBdr,
        line width=1.35pt,
        text width=\fsmIdleW,
        minimum height=\fsmIdleH
    },
    stateServ/.style={
        fsmNode,
        fill=fsmServFill,
        draw=fsmServBdr,
        text width=\fsmServW,
        minimum height=\fsmServH
    },
    stateWarm/.style={
        fsmNode,
        fill=fsmWarmFill,
        draw=fsmWarmBdr,
        line width=1.1pt,
        text width=\fsmAutoW,
        minimum height=\fsmAutoH
    },
    stateHero/.style={
        fsmNode,
        fill=fsmHeroFill,
        draw=fsmHeroBdr,
        line width=1.85pt,
        text width=\fsmAutoW,
        minimum height=\fsmAutoH
    },
    stateStop/.style={
        fsmNode,
        fill=fsmStopFill,
        draw=fsmStopBdr,
        line width=1.1pt,
        text width=3.10cm,
        minimum height=1.05cm
    },
    % --- Estilos de transiciones (flechas canónicas) ---
    transMain/.style={
        -{Stealth[length=2.5mm, width=1.8mm]},
        line width=1.15pt,
        draw=fsmAutoZoneTxt,
        shorten >=1.5pt,
        shorten <=1.0pt
    },
    transWarm/.style={
        -{Stealth[length=2.5mm, width=1.8mm]},
        line width=1.10pt,
        draw=fsmWarmHdr,
        shorten >=1.5pt,
        shorten <=1.0pt
    },
    transServ/.style={
        -{Stealth[length=2.3mm, width=1.7mm]},
        line width=0.95pt,
        draw=fsmServHdr,
        shorten >=1.5pt,
        shorten <=1.0pt
    },
    transRet/.style={
        -{Stealth[length=2.3mm, width=1.7mm]},
        line width=0.90pt,
        draw=fsmIdleBdr,
        shorten >=1.5pt,
        shorten <=1.0pt
    },
    transShort/.style={
        -{Stealth[length=2.0mm, width=1.5mm]},
        line width=1.10pt,
        shorten >=1.5pt,
        shorten <=1.0pt
    },
    transBypass/.style={
        -{Stealth[length=2.5mm, width=1.8mm]},
        line width=1.05pt,
        draw=fsmHeroBdr,
        dashed,
        dash pattern=on 5pt off 3pt,
        shorten >=1.5pt,
        shorten <=1.0pt
    },
    transStop/.style={
        -{Stealth[length=2.5mm, width=1.8mm]},
        line width=1.10pt,
        draw=fsmStopHdr,
        shorten >=1.5pt,
        shorten <=1.0pt
    },
    % --- Etiquetas de transición ---
    elbl/.style={
        font=\small\sffamily,
        text=fsmInk,
        align=center,
        inner sep=1.8pt
    },
    % --- Contenedores y encabezados ---
    zoneServ/.style={
        draw=fsmServZoneBdr,
        dashed,
        line width=0.95pt,
        rounded corners=10pt,
        fill=none,
        inner xsep=0.34cm,
        inner ysep=0.44cm
    },
    zoneAuto/.style={
        draw=fsmAutoZoneBdr,
        dashed,
        line width=1.10pt,
        rounded corners=10pt,
        fill=none,
        inner xsep=0.36cm,
        inner ysep=0.48cm
    },
    zonePill/.style={
        rectangle,
        rounded corners=3.5pt,
        font=\small\bfseries\sffamily,
        inner xsep=7pt,
        inner ysep=2.5pt
    }
]

% ==============================================================================
% 1. COLOCACIÓN RELATIVA DE ESTADOS (ÚNICO ORIGEN EN IDLE)
% ==============================================================================

% --- Núcleo central: IDLE ---
\node[stateIdle] (idle) at (0,0) {{\large\bfseries IDLE}};

% --- Rama Izquierda: Servicio ---
\node[stateServ, left=\fsmColGapL of idle] (char) {{\small\bfseries CHARACTERIZATION}};
\node[stateServ, above=\fsmServGapY of char] (cal) {{\small\bfseries CALIBRATING}};
\node[stateServ, below=\fsmServGapY of char] (pid) {{\small\bfseries MOTOR\_PID}};

% --- Rama Derecha: Ciclo Autónomo ---
\node[stateWarm, right=\fsmColGapR of idle, yshift={0.5*(\fsmAutoH+\fsmAutoGapY)}] (swing)
    {{\normalsize\bfseries SWING\_UP}};
\node[stateHero, right=\fsmAutoGapX of swing] (balance)
    {{\normalsize\bfseries\color{fsmHeroHdr} BALANCE\_PID}};

\coordinate (auto_row1_mid) at ($(swing.south west)!0.5!(balance.south east)$);
\path let \p1 = (swing.west), \p2 = (balance.east),
          \n1 = {\x2 - \x1 - 2*\fsmNodePadX - 1.50pt} in
    node[stateWarm, text width=\n1, anchor=north] (recover)
        at ($(auto_row1_mid)+(0,-\fsmAutoGapY)$) {{\normalsize\bfseries RECOVER\_HOME}};

% ==============================================================================
% 2. CONTENEDORES ELÁSTICOS (FIT) CON ENCABEZADOS TIPO PILL
% ==============================================================================

\coordinate (serv_top_pad) at ($(cal.north)+(0,0.18cm)$);
\node[zoneServ, fit=(cal)(char)(pid)(serv_top_pad)] (serv_zone) {};
\node[zonePill, fill=fsmServZonePill, draw=fsmServZoneBdr, text=fsmServZoneTxt]
    at (serv_zone.north) {Servicio};

\coordinate (auto_top_pad) at ($(swing.north)+(0,0.22cm)$);
\node[zoneAuto, fit=(swing)(balance)(recover)(auto_top_pad)] (auto_zone) {};
\node[zonePill, fill=fsmAutoZonePill, draw=fsmAutoZoneBdr, text=fsmAutoZoneTxt, anchor=west]
    at ([xshift=0.40cm]auto_zone.north west) {Ciclo Autónomo};

% --- Parada Ordenada: STOP_HOME ---
\node[stateStop, below=1.05cm of auto_zone.south] (stophome) {{\normalsize\bfseries STOP\_HOME}};

% ==============================================================================
% 3. TRANSICIONES Y ETIQUETAS ULTRA BREVES (CERO CRUCES)
% ==============================================================================

% --- Entrada inicial (Inicio -> IDLE) ---
\coordinate (init_dot) at ([xshift=-0.70cm, yshift=0.78cm]idle.north);
\fill[fsmIdleHdr] (init_dot) circle (3.0pt);
\draw[transShort, draw=fsmIdleHdr]
    (init_dot) -- ([xshift=-0.70cm]idle.north);
\node[elbl, above=2.5pt, font=\small\bfseries\sffamily, text=fsmIdleHdr]
    at (init_dot) {Inicio};

% --- Entrada global de emergencia (STOP / B1 -> IDLE) ---
\coordinate (stop_b1_dot) at ([xshift=-0.70cm, yshift=-0.78cm]idle.south);
\draw[draw=fsmAlertBdr, fill=white, line width=1.05pt] (stop_b1_dot) circle (3.2pt);
\fill[fsmAlertBdr] (stop_b1_dot) circle (1.5pt);
\draw[transShort, draw=fsmAlertBdr]
    ([yshift=3.2pt]stop_b1_dot) -- ([xshift=-0.70cm]idle.south);
\node[elbl, below=3pt, text=fsmAlertTxt]
    at (stop_b1_dot) {\texttt{STOP / B1}};

% --- Carriles ortogonales en brecha izquierda (Servicio <-> IDLE) ---
\coordinate (laneL_near) at ($(idle.west)+(-\fsmLaneStep,0)$);
\coordinate (laneL_far)  at ($(idle.west)+(-2*\fsmLaneStep,0)$);
\coordinate (gapL_mid)   at ($(serv_zone.east)!0.5!(laneL_far)$);

% IDLE <-> CALIBRATING
\coordinate (cal_in)  at ([yshift=0.22cm]cal.east);
\coordinate (cal_out) at ([yshift=-0.22cm]cal.east);
\draw[transServ, rounded corners=6pt]
    ([yshift=0.78cm]idle.west) -- ([yshift=0.78cm]idle.west -| laneL_near)
    -- (cal_in -| laneL_near)
    -- (cal_in);
\node[elbl, above=2pt] at (cal_in -| gapL_mid) {\texttt{ZERO / B1}};

\draw[transRet, rounded corners=6pt]
    (cal_out) -- (cal_out -| laneL_far)
    -- ([yshift=0.45cm]idle.west -| laneL_far)
    -- ([yshift=0.45cm]idle.west);
\node[elbl, below=2pt, text=fsmMuted] at (cal_out -| gapL_mid) {Listo};

% IDLE <-> CHARACTERIZATION
\draw[transServ]
    ([yshift=0.16cm]idle.west) -- ([yshift=0.16cm]char.east);
\node[elbl, above=2pt] at ([yshift=0.16cm]gapL_mid) {\texttt{CHAR}};

\draw[transRet]
    ([yshift=-0.16cm]char.east) -- ([yshift=-0.16cm]idle.west);
\node[elbl, below=2pt, text=fsmMuted] at ([yshift=-0.16cm]gapL_mid) {Fin};

% IDLE -> MOTOR_PID
\draw[transServ, rounded corners=6pt]
    ([yshift=-0.75cm]idle.west) -- ([yshift=-0.75cm]idle.west -| laneL_near)
    -- (pid.east -| laneL_near)
    -- (pid.east);
\node[elbl, above=2pt] at (pid.east -| gapL_mid) {\texttt{VEL / POS}};

% --- Carril derecho (IDLE -> Ciclo Autónomo) ---
\coordinate (laneR)    at ($(idle.east)+(\fsmLaneStep,0)$);
\coordinate (gapR_mid) at ($(laneR)!0.5!(auto_zone.west)$);

% IDLE -> SWING_UP
\draw[transMain, rounded corners=6pt]
    ([yshift=0.55cm]idle.east) -- ([yshift=0.55cm]idle.east -| laneR)
    -- (swing.west -| laneR)
    -- (swing.west);
\node[elbl, above=2pt] at (swing.west -| gapR_mid) {\texttt{START}};

% IDLE -> RECOVER_HOME
\draw[transWarm, rounded corners=6pt]
    ([yshift=-0.55cm]idle.east) -- ([yshift=-0.55cm]idle.east -| laneR)
    -- (recover.west -| laneR)
    -- (recover.west);
\node[elbl, above=2pt] at (recover.west -| gapR_mid) {\texttt{START} ($\theta \neq 0$)};

% IDLE -> BALANCE_PID (riel superior)
\coordinate (top_rail) at ($(auto_zone.north)+(0,0.58cm)$);
\coordinate (bal_exit) at ([xshift=0.70cm]idle.north);
\draw[transBypass, rounded corners=7pt]
    (bal_exit) -- (bal_exit |- top_rail)
    -- node[pos=0.55, above=2.5pt, elbl, text=fsmHeroHdr] {\texttt{BALANCE}}
       (balance.north |- top_rail)
    -- (balance.north);

% --- Transiciones internas del Ciclo Autónomo ---

% SWING_UP -> BALANCE_PID
\draw[transMain, draw=fsmHeroBdr, line width=1.30pt]
    (swing.east) -- node[midway, above=3pt, elbl, font=\small\bfseries\sffamily, text=fsmHeroHdr] {Captura}
    (balance.west);

% SWING_UP -> RECOVER_HOME (Timeout)
\coordinate (sw_tmo_top) at ([xshift=-0.42cm]swing.south);
\coordinate (sw_tmo_bot) at (sw_tmo_top |- recover.north);
\draw[transWarm]
    (sw_tmo_top) -- node[midway, left=3pt, elbl, text=fsmWarmHdr] {Timeout} (sw_tmo_bot);

% RECOVER_HOME -> SWING_UP (Calma)
\coordinate (sw_rel_top) at ([xshift=0.42cm]swing.south);
\coordinate (sw_rel_bot) at (sw_rel_top |- recover.north);
\draw[transMain, line width=1.15pt]
    (sw_rel_bot) -- node[midway, right=3pt, elbl, text=fsmAutoZoneTxt] {Calma} (sw_rel_top);

% BALANCE_PID -> RECOVER_HOME (Caída)
\coordinate (bal_wd_bot) at (balance.south |- recover.north);
\draw[transWarm, line width=1.15pt]
    (balance.south) -- node[midway, left=4pt, elbl, text=fsmWarmHdr] {Caída} (bal_wd_bot);

% --- Parada ordenada: Ciclo Autónomo -> STOP_HOME -> IDLE ---
\draw[transStop]
    (auto_zone.south) -- node[midway, right=4pt, elbl, text=fsmStopHdr] {\texttt{STOPHOME}} (stophome.north);

\draw[transStop, rounded corners=7pt]
    (stophome.west) -| node[pos=0.25, above=2.5pt, elbl, text=fsmStopHdr] {En $0^\circ$}
    ([xshift=0.70cm]idle.south);

\end{tikzpicture}%
}%
\caption{Máquina de estados finitos (FSM) del firmware de control digital (\texttt{furuta\_control.c}, lazo isócrono \texttt{TIM6} a \qty{1.0}{\kilo\hertz}). Comprende el estado inicial de reposo \texttt{IDLE}, la rama de servicio (\texttt{CALIBRATING}, \texttt{CHARACTERIZATION} y \texttt{MOTOR\_PID}), el ciclo autónomo cerrado (\texttt{SWING\_UP}, \texttt{BALANCE\_PID} y \texttt{RECOVER\_HOME}) y la parada ordenada \texttt{STOP\_HOME}. La transición global de emergencia mediante el comando UART \texttt{STOP} o el botón \texttt{B1} (\texttt{PC13}) fuerza el retorno inmediato a \texttt{IDLE} desde cualquier estado; la descomposición en subfases de \texttt{RECOVER\_HOME} se presenta en la Fig.~\ref{fig:detalle_recover_home}.}
\label{fig:fsm_firmware}
\end{figure*}
